% ============================================================
%  USPSC-Cap3-Metodologia_3.5-Preprocessamento_Textual.tex
%  Seção 3.5 — Pré-processamento Textual
% ============================================================

\section{Pré-processamento Textual}
\label{sec:preprocessamento}

O pré-processamento textual constitui uma etapa crítica em qualquer \textit{pipeline} de PLN, com impacto direto sobre a qualidade das representações vetoriais e, consequentemente, sobre o desempenho do classificador de sentimento \cite{caseli2024cap1}. Para textos provenientes de mídias sociais, essa etapa é particularmente relevante dado o caráter informal, ruidoso e semanticamente denso desses dados, que inclui elementos típicos da linguagem de plataformas digitais --- \textit{hashtags}, menções, URLs e emojis --- que não fazem parte do vocabulário para o qual os modelos de linguagem convencionais foram projetados \cite{santana2024cap34}.

Diferentemente de uma abordagem que aplicaria um único fluxo de limpeza uniforme a todas as abordagens de classificação, este trabalho adota uma \textbf{configuração de pré-processamento específica para cada classificador}, selecionada empiricamente: para cada abordagem, testou-se um conjunto de combinações de etapas de limpeza sobre uma amostra do corpus, comparando a classificação resultante com o rótulo humano e retendo a configuração que produziu a maior taxa de concordância humano-modelo. Essa escolha reconhece que cada abordagem --- léxica ou baseada em modelo de linguagem pré-treinado --- responde de forma distinta ao ruído textual típico do X/Twitter, de modo que a configuração ótima não é necessariamente a mesma entre elas.

As funções de limpeza residem em \textbf{\texttt{app/shared/text\_cleaner.py}}, módulo transversal compartilhado entre os estágios de pré-processamento e classificação, garantindo consistência e evitando duplicação de código. As subseções a seguir descrevem a configuração vencedora de cada abordagem, na ordem em que as etapas são aplicadas.

\subsection{Fluxo de Pré-processamento --- OpLexicon}
\label{subsec:fluxo_oplexicon}

A melhor configuração para o OpLexicon, implementada no método \texttt{preprocess()} da classe \texttt{OpLexiconAnalyzer} (\texttt{app/core/classification/lexicon/op\_lexicon.py}), alcançou \textbf{25,7\% de concordância} humano-modelo e compreende nove etapas sequenciais:

\begin{enumerate}
	\item \textbf{Remoção de URLs}: toda sequência de caracteres iniciada por \texttt{http} ou \texttt{https} é removida do texto, sem substituição por marcador equivalente;
	
	\item \textbf{Remoção de emojis}: os emojis são removidos do texto por meio de \texttt{emoji.replace\_emoji()}, sem substituição por marcador textual ou \textit{name code} equivalente;
	
	\item \textbf{Remoção de menções}: referências a usuários (\texttt{@usuário}) são removidas do texto;
	
	\item \textbf{Remoção do símbolo de \textit{hashtag}}: o caractere \texttt{\#} é removido, mas o termo associado é preservado (e.g., \texttt{\#IBOV} torna-se \texttt{IBOV}) \cite{ranco2015};
	
	\item \textbf{Remoção de pontuação de borda}: caracteres de pontuação de sentença (\texttt{. , ! ?}) são removidos apenas das extremidades de cada \textit{token}, preservando pontuação interna com valor semântico ou estrutural;
	
	\item \textbf{Normalização de espaçamento}: sequências de espaços múltiplos, tabulações e quebras de linha são colapsadas em um único espaço;
	
	\item \textbf{Normalização de caixa com preservação de entidades}: o texto é convertido para letras minúsculas, com exceção de \textit{tickers} de ações (padrão \texttt{[A-Z]\{4,5\}[0-9]\{1,2\}}) e entidades do conjunto \texttt{FINANCIAL\_ENTITIES} (BCB, CVM, IBOV, SELIC, B3, entre outras), que permanecem em maiúsculas;
	
	\item \textbf{Remoção de \textit{stopwords}}: palavras funcionais de alta frequência sem valor semântico para a tarefa são removidas, utilizando a lista de \textit{stopwords} do português da biblioteca \texttt{NLTK} \cite{jurafsky2026};
	
	\item \textbf{Lematização}: as palavras são reduzidas à sua forma canônica (lema) por meio do modelo de língua portuguesa \texttt{pt\_core\_news\_lg} da biblioteca \texttt{spaCy}\footnote{Disponível em: \url{https://spacy.io/}. Acesso em: 12 fev. 2026.}, preservando distinções lexicais relevantes ao domínio (e.g., ``alta'' como substantivo versus ``alto'' como adjetivo). Essa etapa é necessária porque o OpLexicon v3.0 indexa suas entradas por lema, e não por forma flexionada.
\end{enumerate}

\subsection{Fluxo de Pré-processamento --- SentiLex-PT}
\label{subsec:fluxo_sentilex}

A melhor configuração para o SentiLex-PT, implementada no método \texttt{preprocess()} da classe \texttt{SentiLexAnalyzer} (\texttt{app/core/classification/lexicon/senti\_lex.py}), alcançou \textbf{29,5\% de concordância} humano-modelo. Aplica as mesmas oito primeiras etapas descritas para o OpLexicon (\autoref{subsec:fluxo_oplexicon}), na mesma ordem --- remoção de URLs, emojis, menções, símbolo de \textit{hashtag}, pontuação de borda, normalização de espaçamento e de caixa, e remoção de \textit{stopwords} --- \textbf{sem a etapa final de lematização}. Diferentemente do OpLexicon, o SentiLex-PT02 indexa suas entradas tanto por lema quanto por forma flexionada, o que torna a lematização desnecessária para essa abordagem.

\subsection{Fluxo de Pré-processamento --- FinBERT-PT-BR}
\label{subsec:fluxo_finbert}

A configuração vencedora para o FinBERT-PT-BR, implementada no método \texttt{preprocess()} da classe \texttt{FinBertPTBRAnalyzer} (\texttt{app/core/classification/bert/fin bert\_ptbr.py}), alcançou \textbf{40,5\% de concordância} humano-modelo --- a maior entre as abordagens que aplicam alguma etapa de limpeza --- com uma única etapa: \textbf{normalização de espaçamento}, que colapsa sequências de espaços múltiplos, tabulações e quebras de linha em um único espaço.

Esse resultado é notável porque contraria a expectativa inicial de que URLs, emojis, menções e \textit{hashtags} deveriam ser tratados explicitamente antes da inferência. Uma hipótese explicativa é que o tokenizador \textit{WordPiece} interno ao BERT \cite{devlin2019} já é capaz de segmentar esses elementos em subunidades processáveis sem limpeza prévia, e que etapas adicionais de normalização podem remover sinais textuais que, empiricamente, contribuem para a concordância com o julgamento humano nesse corpus específico.

\subsection{Fluxo de Pré-processamento --- BERTimbau Ajustado}
\label{subsec:fluxo_bertimbau}

A configuração vencedora para o BERTimbau ajustado, implementada no método \texttt{preprocess()} da classe \texttt{BERTimbauAnalyzer} (\texttt{app/core/classification/bert/bert\_timbau.py}), alcançou \textbf{89,5\% de concordância} humano-modelo --- a maior entre todas as quatro abordagens --- e compreende sete etapas sequenciais:

\begin{enumerate}
	\item \textbf{Substituição de URLs}: toda sequência iniciada por \texttt{http} ou \texttt{https} é substituída pelo \textit{token} \texttt{[URL]}, preservando o sinal de que um \textit{link} estava presente;
	
	\item \textbf{Substituição de emojis por códigos}: os emojis são convertidos para seus \textit{name codes} por meio de \texttt{emoji.demojize()} (e.g., \includegraphics[width=0.5cm]{USPSC-TA-Textual/png/rotating_light.png} torna-se \texttt{:rotating\_light:}), preservando seu potencial semântico \cite{sprenger2014} de forma processável pelo tokenizador do modelo;
	
	\item \textbf{Substituição de menções}: referências a usuários (\texttt{@usuário}) são substituídas pelo \textit{token} \texttt{[MENTION]};
	
	\item \textbf{Remoção do símbolo de \textit{hashtag}}: o caractere \texttt{\#} é removido, mas o termo associado é preservado;
	
	\item \textbf{Remoção de pontuação de borda}: caracteres de pontuação de sentença são removidos das extremidades de cada \textit{token}, preservando pontuação interna com valor semântico;
	
	\item \textbf{Normalização de espaçamento}: sequências de espaços múltiplos, tabulações e quebras de linha são colapsadas em um único espaço;
	
	\item \textbf{Normalização de caixa com preservação de entidades}: o texto é convertido para minúsculas, com exceção de \textit{tickers} e entidades financeiras, que permanecem em maiúsculas --- consistente com a variante \texttt{cased} do modelo base (\texttt{neuralmind/bert-base-portuguese-cased}), que preserva a distinção entre maiúsculas e minúsculas durante a tokenização \cite{souza2020}.
\end{enumerate}


\subsection{Comparativo entre os Fluxos}
\label{subsec:comparativo_preproc}

O \autoref{tab:comparativo_preproc} sintetiza as diferenças entre as quatro configurações.

\begin{table}[!htbp]
	\caption{Comparativo entre as configurações de pré-processamento vencedoras de cada abordagem}
	\label{tab:comparativo_preproc}
	\begin{tabular}{p{2.6cm}p{2.9cm}p{2.9cm}p{2.3cm}p{2.9cm}}
		\hline
		\textbf{Etapa} & \textbf{OpLexicon} & \textbf{SentiLex-PT} & \textbf{FinBERT-PT-BR} & \textbf{BERTimbau} \\
		\hline
		URLs & Remove & Remove & --- & Substitui por \texttt{[URL]} \\
		\hline
		Emojis & Remove & Remove & --- & Converte p/ \textit{name code} \\
		\hline
		Menções & Remove & Remove & --- & Substitui por \texttt{[MENTION]} \\
		\hline
		Símbolo \texttt{\#} & Remove, mantém termo & Remove, mantém termo & --- & Remove, mantém termo \\
		\hline
		Pontuação de borda & Sim & Sim & --- & Sim \\
		\hline
		Normalização de espaços & Sim & Sim & Sim & Sim \\
		\hline
		Normalização de caixa & Sim, preserva entidades & Sim, preserva entidades & --- & Sim, preserva entidades \\
		\hline
		\textit{Stopwords} & Remove & Remove & --- & --- \\
		\hline
		Lematização & Sim & --- & --- & --- \\
		\hline
		\textbf{Concordância} & \textbf{25,7\%} & \textbf{29,5\%} & \textbf{40,5\%} & \textbf{89,5\%} \\
		\hline
	\end{tabular}
	\legend{Fonte: Elaborado pelo autor com base no código-fonte de \texttt{app/shared/text\_cleaner.py} e dos métodos \texttt{preprocess()} de \texttt{app/core/classification/lexicon/op\_lexicon.py}, \texttt{app/core/classification/lexicon/senti\_lex.py}, \texttt{app/core/classification/bert /finbert\_ptbr.py} e \texttt{app/core/classification/bert/bert\_timbau.py}.}
\end{table}	
	
	
\subsection{Testes Unitários}
\label{subsec:testes_preprocessamento}

O módulo de pré-processamento é acompanhado de um arquivo de testes unitários co-localizado em \texttt{app/shared/text\_cleaner\_tests.py}, executado via \texttt{pytest} a partir da raiz do projeto. Os testes verificam o comportamento de cada função individualmente — cobrindo URLs, emojis, menções, \textit{hashtags}, espaçamento, normalização de caixa, \textit{stopwords} e lematização —, assegurando que alterações futuras no módulo não introduzam regressões silenciosas no fluxo de
pré-processamento.